\section{Implementation}
\label{sec:impl}

The compiler is a header-only C++ library; the CUDA backend is a single header.
A grammar is a C++ type; all lowering happens at compile time, so the generated
scanner is an ordinary templated kernel with no dispatch cost.

\subsection{Front end: shape-free lowering}
\label{sec:frontend}

The front end (\code{include/pars/compile/}) reads only rule definitions:
\begin{enumerate}
  \item \textbf{Recursion analysis.} The rule reference graph is traversed with a
    visited set to compute, per rule, whether it lies on a cycle and how many
    direct self-references it has (cycle glue has none).
  \item \textbf{Automaton.} Each lexical rule body is reified to a regular
    expression in a fixed-capacity arena (no dynamic allocation, so it is
    constant-evaluable by strict front ends such as \textsf{nvcc}) and compiled
    to a DFA by Brzozowski derivatives.  States are canonical derivatives;
    alternatives are normalized by ACI sorting and deduplication, and the
    alphabet is the quotient of $\Sigma$ by membership in every character
    predicate, keeping the DFA tiny even for dense complement predicates.
  \item \textbf{Descriptors.} For each lexical rule we store a \emph{monoid
    descriptor} (\S\ref{sec:backend}).  The unified mask predicate of
    \S\ref{sec:when} decides whether the rule masks at all.
  \item \textbf{Brackets.} A recursive rule is accepted as a well-nested anchor
    only if its top-level definition has exactly the two literal endpoints,
    distinct, with the recursion strictly inside (a conservative, sufficient
    structural criterion); well-nestedness is additionally validated at run time.
  \item \textbf{Diagnostics.} Unsupported recursion or unsupported \code{notp}
    forms are a compile error unless the rule is declared in
    \code{fallback\_rules}.
\end{enumerate}

\subsection{Monoid descriptors}
\label{sec:descriptor}

A descriptor is a small structural value (an NTTP) of one of these kinds:
\begin{itemize}
  \item \textbf{affine}: a delimiter set plus an optional prefix-escape byte
    (executed as $x'=(a\wedge x)\oplus b$);
  \item \textbf{doubled}: a delimiter set (parity on valid input);
  \item \textbf{comment}: a start set and a newline reset;
  \item \textbf{table}: a $\le 8\times 8$ transition table with a
    per-transition mask bit.
\end{itemize}
The kind is read off the DFA; importantly, the execution of \emph{every} kind is
defined by the same automaton table, so a specialization is a performance choice,
not a semantic one.

\subsection{Backends}
\label{sec:backend}

\paragraph{CPU.}
The CPU backend executes descriptors as a per-byte state transform; the running
state is the prefix product (\Rfour).

\paragraph{CUDA.}
The GPU backend packs all descriptors into one 64-bit transform: each descriptor
owns a contiguous range of $4$-bit slots, and we precompute a $256$-entry
per-byte transform table (absolute slot encoding).  Per $32$-byte word we compose
the word transform (each thread), run a \code{thrust} inclusive scan of the
packed transforms, and replay to recover the per-byte state and emit the
per-transition mask.  The scan slot width is \code{total\_states}, and the tables
are staged into shared memory; a compile-time gate rejects
$\code{total\_states}>16$.

\subsection{Schedule}
\label{sec:schedule}

The plan separates compute from movement.  \code{scan\_device} is zero-copy
(device input, device results; only scalar count read-back); \code{scan\_host}
is the only adapter with bulk transfers (pinned, asynchronous), and
\code{scan\_chunk} carries the packed prefix-monoid element and the compaction
offset so a caller can double-buffer chunks across two streams.  The bracket
pairing uses a depth radix sort (\Rsix) rather than a comparison sort.
